\section{Methods}
\label{sec:methods}

\subsection{From a checkpoint to energy, forces and stress}
In molecular dynamics (MD), atomic positions change at every step while the potential's learned parameters stay fixed. LoupMD exports a pretrained MACE checkpoint once, then evaluates energies and their derivatives on the GPU for each new geometry. The evaluator requires a cell periodic in three dimensions. Isolated molecules can be placed in a vacuum box. We describe the CUDA implementation at \texttt{release/v1}, revision \texttt{430c4aa}. Each short-range interaction sums messages from neighbours within the cutoff, coupling sender features to spherical harmonics of the edge direction with distance-dependent weights. Symmetric products of the summed messages build higher-body-order features, following the atomic cluster expansion and MACE~\cite{drautz2019ace,batatia2022mace}. Scalar site readouts, elemental reference energies and any pair repulsion specified by the checkpoint give the short-range energy $E=\sum_i E_i$.

For an edge vector $\mathbf v_e$ from sender to receiver, including its periodic image shift, analytic reverse differentiation gives $\mathbf g_e=\partial E/\partial\mathbf v_e$. The edge contributes $+\mathbf g_e$ to the sender force and $-\mathbf g_e$ to the receiver force. The short-range stress is the sum of the symmetrised outer products $\mathbf v_e\otimes\mathbf g_e$ divided by cell volume.

For a given set of elements, export selects the readout head (e.g.\ \texttt{omat\_pbe} for MACE-MH-1~\cite{batatia2025mh1}) and writes weights, radial tables and kernel signatures to a bundle. The default generic engine covers the MACE-MP materials family~\cite{batatia2025materials}, OMat24- and MatPES-trained checkpoints~\cite{barrosoluque2024omat,kaplan2025matpes}, molecular MACE-OFF~\cite{kovacs2025maceoff}, OMol25-trained MACE-OMOL-0~\cite{levine2025omol} and the local part of MACE-POLAR-1~\cite{batatia2026polar}.

\subsection{Neighbour sums and symmetric products}
On the default path, the first interaction takes a species embedding as its input. For each receiver $i$ and sender species $s$, the kernel sums $b_k(r_{ij})Y_{\ell m}(\widehat{\mathbf v}_{ij})$ over neighbours $j$ of species $s$, using the radial channel $b_k$ for the species pair. A precomputed matrix then applies the embedding, radial output weights and message projection to these sums. The first-layer table stores four unpacked float32 coefficients per channel and interval. At later interactions the sender carries environment-dependent features. Here each radial weight for a convolution path and channel occupies eight bytes per interval, with one exponent shared by its four coefficients. Blocks that use an edge density keep its cubic in a separate float32 table.

DP Compress tabulated Deep Potential embedding networks~\cite{lu2022dpcompress}, and Kov\'acs et al. used cubic splines for MACE radial networks~\cite[\S III.5]{kovacs2025maceoff}. LoupMD fits cubic Hermite tables as in Fu et al.~\cite[App.~B]{fu2026symmetrixxl} and Liu et al.~\cite[\S 3.6]{liu2026foundation}. DPA4C uses quintic Hermite tables~\cite[\S 4.4, Supplement S2.2]{li2026dpa4c}. The exporter fits float64 values and slopes at the knots. It refines the grid by comparing sampled interpolation errors with the radial network's own float32 deviations from float64. Packing can require further refinement. At load time the engine finds the extrema of the difference between the decoded and unpacked cubics and of the corresponding slope difference. In each column, the maximum value error divided by the maximum unpacked value magnitude must be at most $5\times10^{-4}$. The same bound applies to slopes with their own normalisation. The loader checks both ratios over all species pairs and distances, then repeats the checks for $r\geq1$~\AA{}.

The convolution generator takes its nonzero Wigner $3j$ coefficients from \texttt{e3nn}~\cite{geiger2022e3nn}. For on-site products, coupling a lower-order equivariant to a further density feature gives a recursive construction of higher-order features~\cite[Eq.~4]{nigam2020recursive}. LoupMD factors the converted cuEquivariance polynomial descriptor~\cite{nvidia2026cueq} into reusable pair tensors. It selects linearly independent contractions by their incremental operation count and folds the change of basis into the learned weights, which may be shared across species or indexed by element. The reverse kernel reconstructs and differentiates the same contractions. Kov\'acs et al. exploited sparsity in MACE's generalised Clebsch--Gordan products~\cite[\S III.5]{kovacs2025maceoff}. Related generators specialise sparse tensor products in OpenEquivariance~\cite{bharadwaj2025oeq} and MACE operators in Symmetrix-XL~\cite[\S 4]{fu2026symmetrixxl}.

LoupMD streams each receiver's edges, as in Sobek~\cite{chorosajev2026sobek}. Fusing the tensor product and summation avoids storing expanded messages per edge. Kov\'acs et al. fused MACE's one-particle basis with neighbour summation~\cite[\S III.5]{kovacs2025maceoff}, and FlashTP fuses operations within the tensor product~\cite[\S\S 4.2--4.4]{lee2025flashtp}. The reverse pass regenerates radial and angular terms, as in Symmetrix-XL~\cite[\S 3]{fu2026symmetrixxl} and DPA4C's single interaction~\cite[\S 4.4]{li2026dpa4c}. The optional OMAT-0 medium implementation uses Fuchs's receiver-streaming schedule for channel-preserving tensor products from OpenEquivariance pull request~223~\cite{fuchs2026uvupr}. Its first interaction evaluates piecewise Chebyshev fits in the power basis~\cite{trefethen2013approx}.

\subsection{Linear maps and arithmetic}
On the whole-cell path, successive linear consumers allow the engine to omit an intermediate atom state and its gradient. For example, a product feature $y$ projected by $W_p$ and then by the next layer's input map $W_u$ contributes $y(W_pW_u)$. The loader composes the matrices in float64 and stores the result in float32. This extends through skip connections for supported short-range models when intermediate readouts are linear or absent and the last skip can be composed with the final readout. For OMOL-0 and MH-1, species-indexed tables can also supply first-layer residual and skip terms. Liu et al. use a species-block lookup for the one-hot self-connection~\cite[\S 3.6]{liu2026foundation}.

On first use of an eligible matrix shape, the engine times candidate kernels with CUDA events outside graph capture and caches the choice. The candidates include cuBLAS, custom single-instruction, multiple-thread (SIMT) kernels and, for short-range models, a CUTLASS-based tensor-core hybrid~\cite{thakkar2023cutlass}. Error-corrected products have precedent in \texttt{cuda\_mace}~\cite[\S III.5]{kovacs2025maceoff} and compensated tensor-core arithmetic~\cite{ootomo2022recovery}. Accumulator behaviour depends on the hardware~\cite{fasi2021tensorcores}. LoupMD truncates each float32 operand $x$ to an 11-bit significand $x_h$ and scales its residual as $x_l'=2^{11}(x-x_h)$. One TensorFloat-32 (TF32) matrix instruction multiplies the high parts. A single FP16-input matrix instruction accumulates both high--low cross products in FP32. The kernel rescales the correction by $2^{-11}$, adds it in FP32 and omits the low--low term. The FP16 conversion cannot overflow if every element of both input matrices satisfies $|x|<32\,752$. If a standalone whole-cell evaluation returns a non-finite energy, force or requested stress with the hybrid loaded, the engine disables it and retries once. A non-finite final energy in an MD block requires a restart from saved state. POLAR uses float32 SIMT matrix products by default.

\subsection{Tiling the forward and reverse passes}
If a supported short-range system exceeds the whole-cell memory budget, spatial tiles gather the receiver and neighbouring sender rows for each layer. With two interactions, plain tiling first writes every atom's first-layer output. A second sweep evaluates and differentiates the last layer by tile, accumulating the gradient of each first-layer row it reads. A third sweep differentiates the first layer with those gradients. Fu et al.~\cite[\S 5]{fu2026symmetrixxl} and Liu et al.~\cite[\S 2.1]{liu2026foundation} describe this three-sweep dependency in MACE. Normally the third sweep repeats the first-layer forward calculation. For eligible density blocks with a linear readout, plain tiling skips this replay if memory can hold the projected messages and density.

For two interactions, the wavefront shortens the lifetime of the connecting rows. It computes a tile's first-layer rows before their first second-layer consumer, then replays and differentiates the first layer once the last consumer has contributed its gradient. A pool holds features and gradients for the maximum number of simultaneously live rows. For example, a recorded OMAT-0 medium run on an RTX 5060 Ti held 119\,833 of 1\,081\,600 first-layer rows (11.1\%) for a SrTiO$_3$ rod with 16 tile planes.\footnote{LoupMD profiling record, \path{STATUS_AND_PROFILE.md}, \S 6.22, 16 November 2026.} Temporary convolution storage scales with the active tile. Positions, forces and neighbour lists still occupy space for the full system. Where the image shift is representable, each tile edge stores a sender and an integer shift in eight bytes, with the receiver implicit in its row. Thaler and Keller describe a different list compression based on Hilbert-ordered clusters~\cite{thaler2026compressedneighbors}. The planner can place lists or inter-layer state in pinned host memory. Deeper short-range models use plain tiling.

\subsection{Neighbour lists and molecular dynamics}
GPU cell bins supply a compressed sparse row list of neighbours within the force cutoff plus a skin $s$. In the lane-group implementation, several GPU lanes cooperate to count and fill one receiver row. Each force pass selects neighbours within the physical cutoff. At fixed cell, the list remains valid while every atom stays within $s/2$ of its build position, following the buffered-list construction of Verlet~\cite{verlet1967fluids}. GPU neighbour searches also underlie HOOMD~\cite{anderson2008gpumd}, GROMACS~\cite{pall2020gromacs} and TorchMD-Net~\cite{pelaez2024torchmdnet}.

Whole-cell, fixed-cell MD advances positions and checks the skin on the device. A CUDA Graph switch~\cite{nvidia2026cudagraphs} selects the force pass and second integration stage or, where supported, a neighbour rebuild. If the edge count or longest row exceeds the captured capacity, the host finishes the rebuild and recaptures the graph. A separate guard suppresses the force pass if a rebuilt row violates ascending sender-species order. The host reports that error on the next status publication or at block end. Graph capture has also been used for regular GROMACS steps~\cite{gray2023gromacsgraphs}, TorchMD-Net evaluations~\cite{pelaez2024torchmdnet} and fixed-capacity MACE chunk passes~\cite[\S 3.6]{liu2026foundation}.

Fixed-cell integration follows the Atomic Simulation Environment (ASE) implementation of velocity Verlet~\cite{larsen2017ase}, Langevin integration~\cite{vandeneijnden2006langevin}, stochastic velocity rescaling~\cite{bussi2007csvr}, Nos\'e--Hoover chains~\cite{martyna1992nhc} and Berendsen equilibration~\cite{berendsen1984bath}. Experimental pressure control for whole-cell, short-range models includes BAOAB splitting~\cite{leimkuhler2013baoab} with a Langevin barostat~\cite{quigley2004langevin} and Martyna--Tobias--Klein dynamics~\cite{martyna1994mtk,tuckerman2006liouville}. Relaxation follows ASE's L-BFGS~\cite{larsen2017ase,liu1989lbfgs} or FIRE~\cite{bitzek2006fire} updates. Eligible fixed-cell L-BFGS runs keep the displacement and gradient history on the GPU.

The standard \texttt{loupmd bench} protocol runs Langevin dynamics at 300~K with a 1~fs step, friction $0.01$~fs$^{-1}$ and a 0.5~\AA{} skin. It reports the median of six block times per step, each from 256 moving-atom steps timed with CUDA events in the engine. Export, initial tuning and 100 warm-up steps precede the measurement.

\subsection{Reciprocal-space terms}
MACE-POLAR-1 updates atomic charge distributions using local features and the potential obtained from a direct reciprocal-space sum over the whole cell~\cite{batatia2026polar}. For stress, LoupMD includes the change of reciprocal geometry under strain, following the upstream MACE electrostatic-stress correction~\cite{baldwin2026polarstress}. For suitable short-range checkpoints, the user can add coordination-dependent DFT-D3 dispersion~\cite{grimme2010d3}, including Becke--Johnson damping~\cite{grimme2011bj}. In three-dimensionally periodic cells, the optional split sum applies the FourierD3 factorisation~\cite[\S\S II.3--II.4]{valeeva2026fourierd3} to the reference $C_6$ matrix indexed by species and coordination reference. Projecting each atom's coordination weights onto its species' eigenvector components yields signed channels for $C_6$. LoupMD retains modes with eigenvalues $|\lambda_k|>10^{-12}\max_j|\lambda_j|$ and rescales the same channels by species for $C_8$. The damped pair energy minus its smooth part is summed to a real-space cutoff. Cardinal B-spline interpolation on a grid and fast Fourier transforms sum the smooth part over periodic images, following smooth particle-mesh Ewald~\cite{essmann1995spme} and dispersion Ewald~\cite{iseleholder2012dispersion}. Coordination has its own cutoff, and its derivatives enter forces and stress.

\subsection{Validation against the checkpoint}
For each tested structure and kernel variant, the development gate compares the engine error $d_q$ with the difference $d_q^{32}$ between float32 and float64 PyTorch MACE. For energy, both errors are absolute differences per atom. Forces and stresses use the maximum absolute component difference. All three properties must satisfy
\begin{equation}
d_q \leq \max(3d_q^{32},\epsilon_q), \qquad q\in\{E,F,\sigma\},
\end{equation}
with floors of $0.01$~meV/atom, $10^{-5}$~eV/\AA{} and $10^{-7}$~eV/\AA$^3$, respectively. The float64 MACE calculator provides the energy, force and short-range stress targets. For POLAR stress the development test instead differentiates the native float64 energy with respect to strain, including the dependence of positions, periodic shifts, cell, reciprocal cell and volume. Dividing the symmetrised derivative by volume gives $\sigma^{64}_{\rm strain}$. The candidate error and the precision allowance then use different targets,
\[
 d_\sigma=\max_{ab}\big|\sigma^{\rm Loup}_{ab}-(\sigma^{64}_{\rm strain})_{ab}\big|,\qquad
 d_\sigma^{32}=\max_{ab}\big|(\sigma^{32}_{\rm calc})_{ab}-(\sigma^{64}_{\rm calc})_{ab}\big|,
\]
where $\sigma_{\rm calc}$ denotes the ordinary MACE calculator output at the indicated precision.

The development command \verb|python tools/dev/test_gate.py MODEL --set crystals| takes the POLAR checkpoint path as \texttt{MODEL}. The three r$^2$SCAN input structures are NaCl (mp-22862), NaNbO$_3$ (mp-3671) and NbCl$_5$ (mp-23307) from the Materials Project.\footnote{The development script looks for \texttt{NaCl.cif}, \texttt{NaNbO3.cif} and \texttt{NbCl5.cif} in \path{ResearchModule/Structures/Benchmark/R2SCAN} under the directory named by \texttt{LOUP\_DATA\_ROOT}.} The test replicates the cells, strains them and perturbs atomic positions using a fixed seed. It excludes \texttt{NaNbO3\_c0} by default for POLAR-1-S because native float32 MACE is unstable on that configuration. The organic test set uses molecular crystals and water. At this revision the packaged \texttt{loupmd validate} command uses centred differences of float64 energies under strain with step $10^{-5}$ as its POLAR stress target.
